\section{Related literature}
\label{sec:literature}

The model sits at the intersection of three literatures: Bayesian VAR estimation for forecasting, set identification of structural shocks by sign and zero restrictions, and the use of empirical suites inside central banks. This section situates the replication in each.

\subsection{Bayesian VARs for macroeconomic forecasting}

The VAR entered empirical macroeconomics with \citet{sims1980} as a deliberately atheoretical alternative to large simultaneous-equation models, but unrestricted VARs are profligate with parameters: with $k=8$ variables and $p=4$ lags the system of Section~\ref{sec:model} has $8\times(32+1+6)=312$ conditional-mean coefficients against 126 quarterly observations. The Minnesota prior of \citet{doan1984} and \citet{litterman1986} made the VAR usable at this scale by shrinking each equation towards an independent random walk, with tightness decaying in the lag order. \citet{kadiyala1997} systematised the conjugate normal-inverse-Wishart (NIW) implementation that preserves a closed-form posterior, and \citet{sims1998} developed the general dummy-observation technology --- including the sum-of-coefficients and single-unit-root priors --- through which such priors are imposed by augmenting the data. \citet{banbura2010} showed that with shrinkage scaled appropriately in the cross-section, Bayesian VARs remain competitive forecasters even with dozens or hundreds of variables, a result that underpins the modern central-bank practice of running medium and large BVARs alongside structural models. \citet{giannone2015} closed the loop by treating the prior hyperparameters themselves as parameters to be chosen by maximising the closed-form marginal likelihood --- the approach the replication implements as an option (Section~\ref{sec:estimation}) alongside fixed default tightness. The Covid-19 observations pose a well-known outlier problem for all of these systems; the replication follows the dummy-variable spirit of the pandemic prior of \citet{cascaldigarcia2022}, absorbing 2020Q1--2021Q2 with one exogenous dummy per quarter.

\subsection{Sign and zero restrictions in structural VARs}

Identification by sign restrictions descends from \citet{faust1998}, \citet{canova2002} and, canonically, \citet{uhlig2005}, who replaced zero-style exclusion restrictions with inequality restrictions on impulse responses and accepted that the structural impact matrix is then only \emph{set}-identified: many rotations $Q$ of the Cholesky factor are consistent with the data and the restrictions. \citet{rubioramirez2010} provided the general theory of identification for SVARs together with efficient algorithms for drawing orthogonal matrices from the Haar measure. \citet{arias2018} is the immediate methodological source for this model: it gives an exact algorithm for drawing from the posterior of structural parameters under \emph{combined} zero and sign restrictions, constructing each column of $Q$ in the null space of its zero restrictions and correcting, via importance weights, for the fact that the recursive construction does not sample uniformly from the restricted set. Section~\ref{sec:identification-math} states the algorithm as implemented, and Section~\ref{sec:implementation} reports the weights' effective sample size.

The sign-restriction approach has attracted well-known critiques. \citet{fry2011} emphasise that reporting pointwise medians across accepted draws mixes different structural models (``median is not a model''); the replication follows the source paper in reporting pointwise medians with 68 and 90 per cent bands, so this caveat applies here exactly as it applies to the Bank's own outputs. \citet{baumeister2015} show that the seemingly agnostic uniform (Haar) prior over rotations is informative about impulse responses, and argue for explicit priors on structural parameters; the importance-weighting step of \citet{arias2018}, which the replication implements, addresses the related but distinct problem of sampling uniformly \emph{conditional} on the zero restrictions, not the deeper prior-sensitivity point. \citet{antolindiaz2018} add narrative restrictions on the signs of specific historical shocks --- a natural extension for this model given the clear 2008 and 2022 episodes in its estimated shock series --- and \citet{chan2025} develop the permutation-invariance results that the source paper and the replication use to reduce rejections without distorting the posterior. \citet{kilian2017} provide the standard book-length treatment of the identification, inference and decomposition machinery used throughout Section~\ref{sec:model}.

\subsection{VARs in central-bank forecasting suites}

Central banks rarely rely on a single model. The Bank of England's forecasting platform is organised around a central organising model --- COMPASS, a New Keynesian DSGE --- surrounded by a large ``suite'' of statistical models used for data-consistent forecasting, cross-checking and risk assessment, of which Bayesian VARs are a prominent part \citep{burgess2013}. The SVAR replicated here is a modern member of that ecosystem: as \citet{brignone2025} describe, it is run every forecast round to give the Monetary Policy Committee a \emph{structural} reading of incoming data --- which shocks moved the forecast --- rather than a policy simulation. Its companion methodology paper \citep{brignone2026} generalises the forecast-revision decomposition of Section~\ref{sec:forecasting-math}. The conditional-forecasting tradition of \citet{waggoner1999}, in which VAR forecasts are conditioned on assumed paths for some variables (market interest-rate paths, energy futures), is the natural next step for the replication and is discussed as an extension in Section~\ref{sec:limitations}; the machinery of Section~\ref{sec:forecasting-math} already contains the required moving-average representation. Within PolicyEngine Macro the SVAR plays the same role in miniature: it is the baseline and conditioning member of an open suite \citep{policyengine}, alongside reform-scoring models, mirroring at open-source scale the division of labour that \citet{burgess2013} describe inside the Bank.
